\section{A restricted CausalML route for heterogeneous offer effects}
\label{exec14:causalml}

Flexible response prediction becomes useful to a bank only when its contrast
matches the offer, population and information available at the decision date.
The inspected Uber CausalML revision declares version 0.17.0, Python at least
3.11, eighteen required dependencies and an Apache-2.0 root notice. Its
T-learner and R-learner provide plausible reusable components for a static
offer comparison. This is a source-based selection, without an installed
environment or measured interval performance \citep{exec14_causalml_primary}.

The T-learner fits separate response models for each treatment and the shared
control. For a binary offer $A$, write
$m_a(x)=E(Y\mid A=a,X=x)$. Under consistency, conditional exchangeability,
overlap and a suitable sampling law, the conditional mean contrast is
$\tau(x)=m_1(x)-m_0(x)$. It is a contrast of conditional means, not an identified
individual customer's two potential outcomes. The target population then
determines how these contrasts should be averaged.

An exact stipulated account population makes the difference visible. Two
pre-offer strata have equal population mass and offer probabilities $1/5$
and $4/5$. Their mean recoveries without an offer are USD~100 and USD~120;
with an offer, USD~105 and USD~135. Suppose exchangeability holds within each
stratum. The conditional effects are 5 and 15, hence
\[
\operatorname{ATE}=\tfrac12(5)+\tfrac12(15)=10,\qquad
\operatorname{ATT}=\tfrac15(5)+\tfrac45(15)=13,\qquad
\operatorname{ATC}=\tfrac45(5)+\tfrac15(15)=7.
\]
The observed treated mean is 129 and the untreated mean is 104, giving an
unadjusted difference of 25. All four numbers are correct for their stated
calculations, but only one is the population ATE. The adapter must retain
the population measure instead of treating the presence of a number named
``effect'' as sufficient \citep{exec14_causalml_oracle}.

If the offer costs USD~8 and the imposed means are known, offering only to the
second stratum gives expected gain $\tfrac12(15-8)=3.5$ dollars per eligible
account, compared with 2 dollars for offering to everyone. This is a planning
calculation with known means. Selecting a rule from fitted effects introduces
uncertainty and selection; these dollar amounts do not establish a learned
policy's value or observed bank ROI.

The inspected implementation returns contrast columns in its treatment-label
order. The wrapper should use that returned mapping, preserve account and
feature IDs, and state treatment-minus-control units. The initial route should
remain binary: in the inspected multi-treatment path, the ATE point average
uses prediction rows while the variance terms use each treatment/control
subset. The populations therefore require a separate analysis before a
multi-arm uncertainty contract is promised \citep{exec14_causalml_primary}.

Uncertainty also has several meanings. The selected CATE interval path uses
prediction percentiles from refitted samples; the default ATE expression
combines arm residual variances and fitted-effect variation. The base bootstrap
resamples rows. None of these operations automatically accounts for household
clusters, repeated account histories or common calendar shocks. With a treated
fraction of $1/100$ and a resample of size 100, the probability of selecting
no treated row is $(99/100)^{100}$, about 0.366. A failed arm cannot be silently
discarded to improve reported interval availability. Changing to stratified
or block resampling changes the procedure and needs its own declared target
and analysis \citep{exec14_causalml_primary,exec14_causalml_oracle}.

The R-learner offers a related computational alternative. Its inspected caller
forms outcome residuals and a weighted regression with pseudo-outcome
$(Y-\widehat m(X))/(A-\widehat e(X))$ and weight
$(A-\widehat e(X))^2$. In the deterministic conditional-mean example above,
substituting the true $m$ and $e$ makes that pseudo-outcome equal to the stratum
effect for both treatment values. The identity explains the intended
orthogonalization; it does not validate learned nuisance functions, delegated
propensity folds or dependence-aware cross-fitting. Those delegated paths
must be reviewed before a broader route is enabled.

Observed uplift curves are another distinct object. The inspected metric
uses true effects if supplied and otherwise uses treated/control observed
means with interpolation. A curve can be computable under confounding and
still fail to identify policy value. An ordinary exact stratified estimator
therefore remains a useful comparator, alongside separately reviewed EconML
and grf paths.

The preferred proposal is an optional pinned T-learner service with explicit
base regressor, binary arms and independent whole accounts. Preserve every
seed, environment identity, population weight, pre-offer feature clock and
outcome horizon. Reject missing arms, unsupported sampling, stale row mappings,
nonfinite output and unsupported targets with a reason. Keep time-varying
treatment, latent-state inference and multi-arm uncertainty outside this first
contract. Full adoption would import several scientific contracts at once;
porting would assume their maintenance without a demonstrated need. Executable
admission remains deferred until input/import paths, dependency notices and
target, column, interval and failure fixtures have been reviewed and run.
\label{exec14:causalml-end}
